\section{数据从哪里来：Web、Code、Curated 与 Synthetic}

有了 token budget 后，第二个问题是 source。不同来源不仅内容不同，还带来不同的抓取、许可、去重、格式和 verification 成本。课程把来源分成 web data、GitHub/code、curated sources 和 synthetic data；StarCoder 则展示 code corpus 怎样做成可检查的数据产品。

\subsection{Source taxonomy：来源决定后续治理成本}

上一章只估算 token budget，本节开始回答这些 tokens 从哪里来，因为 source 会决定后续的 filter、license 和 evaluation 成本。Web 覆盖广但噪声大，GitHub 有结构和许可元数据但含重复与 secrets，curated sources 质量高但规模有限，synthetic data 可定向生成却依赖 teacher、prompt 与 seed diversity。

\lecturefigure{slide-26.jpg}{LLM training data 的四类来源}{Loubna Ben Allal 官方 deck，第 26 页}

\begin{knowledgebox}{来源不是互斥标签}
代码仓库包含 README/Markdown，自然语言网页包含代码，synthetic prompts 又常从 web/curated documents 派生。实际 mixture 应按 provenance 和 transformation lineage 记录，而不是只写最终 domain 百分比。
\end{knowledgebox}

\subsection{Common Crawl：规模大到必须先过滤再下载/处理}

Common Crawl 提供周期性网页快照，raw corpus 可达数百 TB。直接 tokenize 不仅昂贵，还会保留 boilerplate、spam、machine-generated pages、duplicates、PII 与低质量文本，因此通常复用已有 filtered web dataset 或做 distributed filtering。

\lecturefigure{slide-27.jpg}{Web data：Common Crawl 的规模与过滤负担}{Loubna Ben Allal 官方 deck，第 27 页}

\begin{warningbox}{抓到网页不等于拥有训练许可}
可访问性、robots、copyright、terms、personal data 与 jurisdiction 是不同问题。技术 pipeline 能记录来源和移除样本，却不能自动解决所有法律与伦理判断。
\end{warningbox}

\teachervoice{00:13:20--00:16:20，讲者强调 Common Crawl 的原始规模和 filtering cost，建议理解已有 dataset 的处理过程，而不是把 raw crawl 当作即用数据。}

\subsection{FineWeb：更好的 filter 用 early-training curve 证明}

FineWeb slide 比较多个 web corpora 上的小模型训练曲线。读图时应先确认相同 architecture、token budget、evaluation 和 seed，再看哪个 corpus 更快达到较高 score。曲线支持的是 filter pipeline 在该设置下优于对照，不是任何模型/任务都最优。

\lecturefigure{slide-29.jpg}{FineWeb ablation：不同 web corpora 的 early-training performance}{Loubna Ben Allal 官方 deck，第 29 页}

\begin{importantbox}{Data quality 的 operational definition}
高质量不是“看起来像好文章”。在预训练里，它应表现为更低 validation loss、更快 downstream learning、更少 memorization/toxicity/PII、更好 domain coverage，且在 matched compute 和多 seeds 下成立。
\end{importantbox}

\subsection{The Stack：代码数据需要 raw collection、inspection 与 opt-out}

代码不像纯文本网页：文件有 repository、path、license、language、stars、commit history 等结构，也可能包含 generated files、vendor copies、secrets 与 benchmark solutions。The Stack pipeline 因此不仅抓取，还要语言检测、license/metadata 管理、dedup 和社区 inspection。

\lecturefigure{slide-30.jpg}{Code data：The Stack 的 collection 与 processing 概览}{Loubna Ben Allal 官方 deck，第 30 页}

\lecturefigure{slide-31.jpg}{The Stack v1：GitHub collection、raw data、dedup、inspection 与 opt-out}{Loubna Ben Allal 官方 deck，第 31 页}

\begin{knowledgebox}{Opt-out 在 pipeline 中的位置}
Opt-out 需要 stable repository/file identifiers、公开查询工具、处理请求的版本化列表，以及 future retraining 时的排除机制。它不是训练完成后删一个网页链接，而是 data lineage 的一部分。
\end{knowledgebox}

\teachervoice{00:17:10--00:19:10，讲者说明 BigCode 在 v1 发布 inspection tool 和 opt-out form，让开发者能检查代码并请求未来训练排除；这是开放数据项目的核心设计。}

\subsection{The Stack v2：Software Heritage 与更多开发过程数据}

v2 不再主要依赖直接 clone GitHub，而从 Software Heritage archive 提取，并加入 issues、pull requests、commits、Jupyter/Kaggle notebooks、math/coding datasets。这样 corpus 不只教最终代码，还包含问题、修改、讨论和执行环境。

\lecturefigure{slide-32.jpg}{The Stack v2：Software Heritage、过滤与多源开发数据}{Loubna Ben Allal 官方 deck，第 32 页}

\lecturefigure{slide-33.jpg}{The Stack v1 与 v2：raw size、filtered size、languages 与 tokens}{Loubna Ben Allal 官方 deck，第 33 页}

\begin{importantbox}{Raw scale 与 usable scale 的差距就是工程量}
v2 raw size 约为 v1 的十倍，而过滤后仍约为 v1 的四到五倍。被删除的绝大部分不是“浪费”，而是 duplicates、低质量、unsupported formats、PII 或不合适语言；过滤率本身不代表好坏，需看保留 diversity 与模型结果。
\end{importantbox}

\teachervoice{00:19:10--00:20:20，讲者把 v2 的关键改进放在 archive source、更多语言和辅助开发数据，而不仅是“token 更多”。}

\subsection{Synthetic data：从 Phi 到可控生成 corpus}

Textbooks Are All You Need/Phi 系列展示用 GPT-3.5/4 生成高质量教材式数据，可以用远少于 web-scale 的 corpus 训练强小模型。它引发的不是“synthetic 取代真实数据”定论，而是把 generator、prompt、seed、verification 变成新数据工程变量。

\lecturefigure{slide-34.jpg}{Synthetic data 的催化案例：Textbooks Are All You Need}{Loubna Ben Allal 官方 deck，第 34 页}

设 seed/context 为 $s$、prompt template 为 $p$、generator 为 $G$，synthetic sample 为
\[
y\sim G(\cdot\mid p,s).
\]
数据质量取决于 $s$ 的覆盖、$p$ 的约束、$G$ 的能力与后验 verifier $v(y)$，不能只靠采样数量。

\begin{warningbox}{Synthetic feedback loop}
若 generator 和 verifier 共享盲区，重复生成会放大风格同质化、事实错误与 representation collapse。应记录 teacher/version、seed provenance、rejection reasons、novelty 与 human spot checks。
\end{warningbox}

\subsection{Cosmopedia：用 web/curated seeds 扩大多样性}

Cosmopedia 使用 Mixtral-8x7B 生成约 25B tokens 的教材、博客、故事等。核心技巧不是一句“让模型写课本”，而是从 web samples、Stanford courses、WikiHow 等抽取 context，把 topic 与具体 seed 一起放入 prompt，限制 hallucination 并增加主题分散度。

\lecturefigure{slide-35.jpg}{Cosmopedia：开放模型生成的大规模 synthetic text dataset}{Loubna Ben Allal 官方 deck，第 35 页}

\lecturefigure{slide-36.jpg}{Cosmopedia 的数据组成：web seeds、curated sources 与多种体裁}{Loubna Ben Allal 官方 deck，第 36 页}

\begin{knowledgebox}{Synthetic corpus 的四项审计}
\begin{enumerate}
  \item \textbf{Coverage：}topic、language、difficulty、style 是否分散？
  \item \textbf{Grounding：}生成是否受 seed/context 约束？
  \item \textbf{Verification：}事实、代码、数学是否有自动/人工检查？
  \item \textbf{Lineage：}teacher、prompt、seed、filter 与 license 是否可追踪？
\end{enumerate}
\end{knowledgebox}

\teachervoice{00:20:20--00:24:20，讲者说 synthetic data 很新、很有前景，也“very tricky”获得 diversity。Cosmopedia 用大量 web/curated seeds，而不是从少数空泛 prompts 无限扩写。}

\subsection{本章小结}

数据 source 决定后续处理成本。Web 提供覆盖，code 提供结构，curated sources 提供高 signal，synthetic data 提供定向扩展；任何来源都必须附带 provenance、filter、license/opt-out 与 evaluation 设计。

